% Curriculum vitae of Kaili Cao: the full version (Kaili_Cao_CV.tex) and a short one (Kaili_Cao_CV_short.tex)
% both input this file. Build both with `latexmk` in this folder. \full{...} marks what the short version leaves out.
% One column, standard headings and a Unicode text layer, so that automatic parsers read it (see kccv.sty).
\documentclass[12pt]{article}
\ifdefined\shortversion
  \newcommand{\full}[1]{}
  \newcommand{\cvversionlabel}{Curriculum Vitae (short version; full version at kailicao.github.io)}
\else
  \newcommand{\full}[1]{#1}
  \newcommand{\cvversionlabel}{Curriculum Vitae}
\fi
\usepackage{kccv}
\hypersetup{pdftitle={Kaili Cao: Curriculum Vitae},pdfauthor={Kaili Cao},pdfsubject={Curriculum vitae},
  pdfkeywords={astrophysics, cosmology, weak gravitational lensing, image processing, stellar physics, asteroseismology, AI in science}}

\ifdefined\shortversion
  \setlength{\cvitemsep}{3pt}
  \titlespacing*{\section}{0pt}{12pt}{6pt}
\fi
\begin{document}

\cvname{Kaili Cao}
\cvtagline{Eric and Wendy Schmidt AI in Science Postdoctoral Fellow, Michigan Institute for Data and AI in Society (MIDAS), University~of~Michigan, Ann Arbor, Michigan, United States}
Email: \href{mailto:kailicao@umich.edu}{kailicao@umich.edu}; Website: \href{https://kailicao.github.io/}{kailicao.github.io}; ORCID: \href{https://orcid.org/0000-0002-1699-6944}{0000-0002-1699-6944}\par
GitHub: \href{https://github.com/kailicao}{github.com/kailicao}; LinkedIn: \href{https://www.linkedin.com/in/kaili-cao-1840b5229/}{linkedin.com/in/kaili-cao-1840b5229}\par

\section{Positions}

\cventry{2026 -- 2028}{Schmidt AI in Science Postdoctoral Fellow}{University~of~Michigan}{Ann Arbor}{United States}{}
\cventry{2024 -- present}{Member}{Asteroseismology Using The Galactic Bulge Time Domain Survey (Roman Space Telescope Wide-Field Science Team)}{Columbus}{United States}{}
\cventry{2023 -- present}{Member}{\Colorhref{https://roman-hlis-cosmology.caltech.edu/}{Maximizing Cosmological Science with the Roman High Latitude Imaging Survey} (Roman Space Telescope Project Infrastructure Team)}{Pasadena}{United States}{}
\cventry{2023 -- 2026}{Graduate Research Assistant}{\Colorhref{https://ccapp.osu.edu/}{Center for Cosmology and AstroParticle Physics (CCAPP)}, The~Ohio~State~University}{Columbus}{United States}{}
\full{\cventry{2019}{Visiting Undergraduate Student}{Massachusetts Institute of Technology}{Cambridge}{United States}{}}

\section{Education}

\cventry{2023 -- 2026}{Doctor of Philosophy, Physics}{The~Ohio~State~University}{Columbus}{United States}
{Dissertation: \textit{Renovating Analysis Infrastructure for the Next Generation of Astronomical Surveys}}
\cventry{2021 -- 2023}{Master of Science, Physics}{The~Ohio~State~University}{Columbus}{United States}{}
\cventry{2020 -- 2021}{31 credits, Computer Science and Technology}{Nanjing~University}{Nanjing}{China}{}
\cventry{2016 -- 2020}{Bachelor of Arts, French Language and Literature}{Nanjing~University}{Nanjing}{China}{}

\section{Honors and Awards}

\cvitem{2026 -- 2028}{\textbf{Eric and Wendy Schmidt AI in Science Postdoctoral Fellowship}, Michigan Institute for Data and AI in Society, University~of~Michigan.}
\cvitem{2023 -- 2026}{Research support as a member of NASA's Roman HLIS Cosmology Project Infrastructure Team and the Roman asteroseismology Wide-Field Science Team (as a PhD student).}

\section{Research Interests}

\cvitem{Cosmology}{Weak Gravitational Lensing, Cosmic Shear Measurement,\newline Parameter Inference, Fisher Forecasts}
\cvitem{Stellar~Physics}{Stellar Models, Stellar Evolution Code, Interior Physics, Asteroseismology, Red Giant Branch}
\cvitem{Methods}{Numerical Simulations, Numerical Methods, Image Processing,\newline Data Analysis, Machine Learning}

\ifdefined\shortversion
\section{Selected Publications}

\cvitem{}{The full list (10 first-author and 14 coauthored papers, each with my contribution) is on \href{https://kailicao.github.io/about/publications.html}{kailicao.github.io/about/publications.html}; ORCID: \href{https://orcid.org/0000-0002-1699-6944}{0000-0002-1699-6944}.}

\subsection{First-Author Publications}

\cvitem{Oct 2026}{\textbf{Kaili Cao}, David H. Weinberg, Vivian Miranda, et al., \Colorhref{https://ui.adsabs.harvard.edu/abs/2026arXiv260100438C}{Fisher Forecasts for Cosmological Yields from $3\!\times\!2$pt Analysis of the Roman Space Telescope High Latitude Imaging Survey},
\textit{The Astrophysical Journal}, Volume 1010, Number 1, October 2026, Page 21, \Colorhref{https://doi.org/10.3847/1538-4357/ae9bc0}{doi:10.3847/1538-4357/ae9bc0}}

\cvitem{Jul 2026}{\textbf{Kaili Cao}, \Colorhref{https://ui.adsabs.harvard.edu/abs/2026arXiv260709862C}{Efficient Optimal Image Reconstruction for the Nancy Grace Roman Space Telescope and Beyond: II. Implementation of \textsc{Effortless}},
\textit{arXiv e-prints}, arXiv:2607.09862 [astro-ph.IM], July 2026, \Colorhref{https://doi.org/10.48550/arXiv.2607.09862}{doi:10.48550/arXiv.2607.09862} [\Colorhref{https://github.com/Roman-HLIS-Cosmology-PIT/effortless}{GitHub}]}

\cvitem{Jul 2026}{\textbf{Kaili Cao}, \Colorhref{https://ui.adsabs.harvard.edu/abs/2026arXiv260706674C}{Efficient Optimal Image Reconstruction for the Nancy Grace Roman Space Telescope and Beyond: I. First Results with \textsc{Effortless}},
\textit{arXiv e-prints}, arXiv:2607.06674 [astro-ph.IM], July 2026, \Colorhref{https://doi.org/10.48550/arXiv.2607.06674}{doi:10.48550/arXiv.2607.06674} [\Colorhref{https://github.com/Roman-HLIS-Cosmology-PIT/effortless}{GitHub}]}

\cvitem{Mar 2026}{\textbf{Kaili Cao} and Roman HLIS Cosmology PIT, \Colorhref{https://ui.adsabs.harvard.edu/abs/2026AJ....171..140C}{Linear Image Regridding and Coaddition with Oversampled Point Spread Functions: Lessons from 1D},
\textit{The Astronomical Journal}, Volume 171, Number 3, March 2026, Page 140, \Colorhref{https://doi.org/10.3847/1538-3881/ae356a}{doi:10.3847/1538-3881/ae356a} [\Colorhref{https://github.com/Roman-HLIS-Cosmology-PIT/effortless}{GitHub}]}

\cvitem{Feb 2026}{\textbf{Kaili Cao}, Christopher M. Hirata, Katherine Laliotis, et al., \Colorhref{https://ui.adsabs.harvard.edu/abs/2026ApJ...998..304C}{Simulating Image Coaddition with the Nancy Grace Roman Space Telescope. IV. Hyperparameter Optimization and Experimental Features},
\textit{The Astrophysical Journal}, Volume 998, Number 2, February 2026, Page 304, \Colorhref{https://doi.org/10.3847/1538-4357/ae3a97}{doi:10.3847/1538-4357/ae3a97} [\Colorhref{https://github.com/Roman-HLIS-Cosmology-PIT/pyimcom.git}{GitHub}]}

\cvitem{Jan 2026}{\textbf{Kaili Cao} and Marc H. Pinsonneault, \Colorhref{https://ui.adsabs.harvard.edu/abs/2026ApJ...997..120C}{Modeling APOKASC-3 Red Giants. I. The First Dredge-up and Red Giant Branch Bump},
\textit{The Astrophysical Journal}, Volume 997, Number 1, January 2026, Page 120, \Colorhref{https://doi.org/10.3847/1538-4357/ae22d1}{doi:10.3847/1538-4357/ae22d1} [\Colorhref{https://github.com/kailicao/mesa_apokasc.git}{GitHub}, \Colorhref{https://zenodo.org/records/15863879}{Zenodo}]}

\cvitem{Dec 2025}{\textbf{Kaili Cao}, \Colorhref{https://ui.adsabs.harvard.edu/abs/2025Math...13.3981C}{ContEvol Formalism: Numerical Methods Based on Hermite Spline Optimization},
\textit{Mathematics}, Volume 13, Issue 24, December 2025, Article Number 3981, \Colorhref{https://doi.org/10.3390/math13243981}{doi:10.3390/math13243981} [\Colorhref{https://github.com/kailicao/ContEvol_formalism.git}{GitHub}]}

\cvitem{Sep 2025}{\textbf{Kaili Cao} and Marc H. Pinsonneault, \Colorhref{https://ui.adsabs.harvard.edu/abs/2025ApJ...990L..26C}{Toward Accurate Asteroseismic Masses for Luminous Giants},
\textit{The Astrophysical Journal Letters}, Volume 990, Number 1, September 2025, Page L26, \Colorhref{https://doi.org/10.3847/2041-8213/adfa12}{doi:10.3847/2041-8213/adfa12}}

\cvitem{Apr 2025}{\textbf{Kaili Cao}, Christopher M. Hirata, Katherine Laliotis, et al., \Colorhref{https://ui.adsabs.harvard.edu/abs/2025ApJS..277...55C}{Simulating Image Coaddition with the Nancy Grace Roman Space Telescope. III. Software Improvements and New Linear Algebra Strategies},
\textit{The Astrophysical Journal Supplement Series}, Volume 277, Number 2, April 2025, Page 55, \Colorhref{https://doi.org/10.3847/1538-4365/adb580}{doi:10.3847/1538-4365/adb580} [\Colorhref{https://github.com/Roman-HLIS-Cosmology-PIT/pyimcom.git}{GitHub}]}

\cvitem{May 2021}{\textbf{Kaili Cao}, David J. Barnes, and Mark Vogelsberger, \Colorhref{https://ui.adsabs.harvard.edu/abs/2021MNRAS.503.3394C}{Studying galaxy cluster morphological metrics with MOCK-X},
\textit{Monthly Notices of the Royal Astronomical Society}, Volume 503, Issue 3, May 2021, Pages 3394--3413, \Colorhref{https://doi.org/10.1093/mnras/stab605}{doi:10.1093/mnras/stab605}}

\else
\section{Publications}

\cvitem{ORCID}{ORCID ID: \Colorhref{https://orcid.org/0000-0002-1699-6944}{0000-0002-1699-6944}; NASA/ADS: \Colorhref{https://ui.adsabs.harvard.edu/search/q=orcid\%3A0000-0002-1699-6944&sort=date+desc}{orcid:0000-0002-1699-6944}}

\subsection{First-Author Publications}

\cvitem{Oct 2026}{\textbf{Kaili Cao}, David H. Weinberg, Vivian Miranda, et al., \Colorhref{https://ui.adsabs.harvard.edu/abs/2026arXiv260100438C}{Fisher Forecasts for Cosmological Yields from $3\!\times\!2$pt Analysis of the Roman Space Telescope High Latitude Imaging Survey},
\textit{The Astrophysical Journal}, Volume 1010, Number 1, October 2026, Page 21, \Colorhref{https://doi.org/10.3847/1538-4357/ae9bc0}{doi:10.3847/1538-4357/ae9bc0}}

\cvitem{Jul 2026}{\textbf{Kaili Cao}, \Colorhref{https://ui.adsabs.harvard.edu/abs/2026arXiv260709862C}{Efficient Optimal Image Reconstruction for the Nancy Grace Roman Space Telescope and Beyond: II. Implementation of \textsc{Effortless}},
\textit{arXiv e-prints}, arXiv:2607.09862 [astro-ph.IM], July 2026, \Colorhref{https://doi.org/10.48550/arXiv.2607.09862}{doi:10.48550/arXiv.2607.09862} [\Colorhref{https://github.com/Roman-HLIS-Cosmology-PIT/effortless}{GitHub}]}

\cvitem{Jul 2026}{\textbf{Kaili Cao}, \Colorhref{https://ui.adsabs.harvard.edu/abs/2026arXiv260706674C}{Efficient Optimal Image Reconstruction for the Nancy Grace Roman Space Telescope and Beyond: I. First Results with \textsc{Effortless}},
\textit{arXiv e-prints}, arXiv:2607.06674 [astro-ph.IM], July 2026, \Colorhref{https://doi.org/10.48550/arXiv.2607.06674}{doi:10.48550/arXiv.2607.06674} [\Colorhref{https://github.com/Roman-HLIS-Cosmology-PIT/effortless}{GitHub}]}

\cvitem{Mar 2026}{\textbf{Kaili Cao} and Roman HLIS Cosmology PIT, \Colorhref{https://ui.adsabs.harvard.edu/abs/2026AJ....171..140C}{Linear Image Regridding and Coaddition with Oversampled Point Spread Functions: Lessons from 1D},
\textit{The Astronomical Journal}, Volume 171, Number 3, March 2026, Page 140, \Colorhref{https://doi.org/10.3847/1538-3881/ae356a}{doi:10.3847/1538-3881/ae356a} [\Colorhref{https://github.com/Roman-HLIS-Cosmology-PIT/effortless}{GitHub}]}

\cvitem{Feb 2026}{\textbf{Kaili Cao}, Christopher M. Hirata, Katherine Laliotis, et al., \Colorhref{https://ui.adsabs.harvard.edu/abs/2026ApJ...998..304C}{Simulating Image Coaddition with the Nancy Grace Roman Space Telescope. IV. Hyperparameter Optimization and Experimental Features},
\textit{The Astrophysical Journal}, Volume 998, Number 2, February 2026, Page 304, \Colorhref{https://doi.org/10.3847/1538-4357/ae3a97}{doi:10.3847/1538-4357/ae3a97} [\Colorhref{https://github.com/Roman-HLIS-Cosmology-PIT/pyimcom.git}{GitHub}]}

\cvitem{Jan 2026}{\textbf{Kaili Cao} and Marc H. Pinsonneault, \Colorhref{https://ui.adsabs.harvard.edu/abs/2026ApJ...997..120C}{Modeling APOKASC-3 Red Giants. I. The First Dredge-up and Red Giant Branch Bump},
\textit{The Astrophysical Journal}, Volume 997, Number 1, January 2026, Page 120, \Colorhref{https://doi.org/10.3847/1538-4357/ae22d1}{doi:10.3847/1538-4357/ae22d1} [\Colorhref{https://github.com/kailicao/mesa_apokasc.git}{GitHub}, \Colorhref{https://zenodo.org/records/15863879}{Zenodo}]}

\cvitem{Dec 2025}{\textbf{Kaili Cao}, \Colorhref{https://ui.adsabs.harvard.edu/abs/2025Math...13.3981C}{ContEvol Formalism: Numerical Methods Based on Hermite Spline Optimization},
\textit{Mathematics}, Volume 13, Issue 24, December 2025, Article Number 3981, \Colorhref{https://doi.org/10.3390/math13243981}{doi:10.3390/math13243981} [\Colorhref{https://github.com/kailicao/ContEvol_formalism.git}{GitHub}]}

\cvitem{Sep 2025}{\textbf{Kaili Cao} and Marc H. Pinsonneault, \Colorhref{https://ui.adsabs.harvard.edu/abs/2025ApJ...990L..26C}{Toward Accurate Asteroseismic Masses for Luminous Giants},
\textit{The Astrophysical Journal Letters}, Volume 990, Number 1, September 2025, Page L26, \Colorhref{https://doi.org/10.3847/2041-8213/adfa12}{doi:10.3847/2041-8213/adfa12}}

\cvitem{Apr 2025}{\textbf{Kaili Cao}, Christopher M. Hirata, Katherine Laliotis, et al., \Colorhref{https://ui.adsabs.harvard.edu/abs/2025ApJS..277...55C}{Simulating Image Coaddition with the Nancy Grace Roman Space Telescope. III. Software Improvements and New Linear Algebra Strategies},
\textit{The Astrophysical Journal Supplement Series}, Volume 277, Number 2, April 2025, Page 55, \Colorhref{https://doi.org/10.3847/1538-4365/adb580}{doi:10.3847/1538-4365/adb580} [\Colorhref{https://github.com/Roman-HLIS-Cosmology-PIT/pyimcom.git}{GitHub}]}

\cvitem{May 2021}{\textbf{Kaili Cao}, David J. Barnes, and Mark Vogelsberger, \Colorhref{https://ui.adsabs.harvard.edu/abs/2021MNRAS.503.3394C}{Studying galaxy cluster morphological metrics with MOCK-X},
\textit{Monthly Notices of the Royal Astronomical Society}, Volume 503, Issue 3, May 2021, Pages 3394--3413, \Colorhref{https://doi.org/10.1093/mnras/stab605}{doi:10.1093/mnras/stab605}}


\subsection{Co-authored Publications in Cosmology}

\cvitem{Jul 2026}{Chun-Hao To, Tae-Hyeon Shin, Chihway Chang, et al. (with \textbf{Kaili Cao} as 10th of 21 authors), \Colorhref{https://ui.adsabs.harvard.edu/abs/2026arXiv260709849T}{Robust Photometry for Roman High-Latitude Imaging Survey Cosmology Using Roman and Rubin Imaging},
\textit{arXiv e-prints}, arXiv:2607.09849 [astro-ph.IM], July 2026, \Colorhref{https://doi.org/10.48550/arXiv.2607.09849}{doi:10.48550/arXiv.2607.09849}}

\cvitem{Jun 2026}{Junzhou Zhang, Chihway Chang, Jiachuan Xu, et al. (with \textbf{Kaili Cao} as 7th of 9 authors), \Colorhref{https://ui.adsabs.harvard.edu/abs/2026arXiv260623781Z}{Accurate modeling for $3\times2$pt analyses in Roman and Rubin: a study of model approximations},
\textit{arXiv e-prints}, arXiv:2606.23781 [astro-ph.CO], June 2026, \Colorhref{https://doi.org/10.48550/arXiv.2606.23781}{doi:10.48550/arXiv.2606.23781}}

\cvitem{Apr 2026}{Katherine Laliotis, Christopher M. Hirata, Emily Macbeth, and \textbf{Kaili Cao}, \Colorhref{https://ui.adsabs.harvard.edu/abs/2026PASJ...78..810L}{Removing correlated noise stripes from the Nancy Grace Roman Space Telescope survey images},
\textit{Publications of the Astronomical Society of Japan}, Volume 78, Issue 2, April 2026, Page 810--823, \Colorhref{https://doi.org/10.1093/pasj/psag020}{doi:10.1093/pasj/psag020} [\Colorhref{https://github.com/Roman-HLIS-Cosmology-PIT/pyimcom.git}{GitHub}]}
% \newline \textit{(My contribution: I led the development of the \textsc{PyImcom} software used as the platform of this work and commented on the manuscript.)}}

\cvitem{Dec 2025}{Openuniverse, Lsst Dark Energy Science Collaboration, Roman Hlis Project Infrastructure, et al. (authors in alphabetical order, including \textbf{K. Cao}), \Colorhref{https://ui.adsabs.harvard.edu/abs/2025MNRAS.544.3799O}{OpenUniverse2024: a shared, simulated view of the sky for the next generation of cosmological surveys},
\textit{Monthly Notices of the Royal Astronomical Society}, Volume 544, Issue 4, December 2025, Pages 3799--3823, \Colorhref{https://doi.org/10.1093/mnras/staf1833}{doi:10.1093/mnras/staf1833}}
% \newline \textit{(``\textbf{K.C.}~led development of the \textsc{PyImcom} software used for IMCOM coadd production and contributed to paper writing.'')}}

\cvitem{Dec 2024}{Katherine Laliotis, Emily Macbeth, Christopher M. Hirata, et al. (with \textbf{Kaili Cao} as 4th of 6 authors), \Colorhref{https://ui.adsabs.harvard.edu/abs/2024PASP..136l4506L}{Analysis of Biasing from Noise from the Nancy Grace Roman Space Telescope: Implications for Weak Lensing},
\textit{Publications of the Astronomical Society of the Pacific}, Volume 136, Number 12, December 2024, Page 124506, \Colorhref{https://doi.org/10.1088/1538-3873/ad9bec}{doi:10.1088/1538-3873/ad9bec}}
% \newline \textit{(My contribution: I led the development of the \textsc{PyImcom} software used for image combination, participated in discussions about data analysis methods, and commented on the manuscript.)}}

\cvitem{Mar 2024}{Masaya Yamamoto, Katherine Laliotis, Emily Macbeth, et al. (with \textbf{Kaili Cao} as 7th of 19 authors), \Colorhref{https://ui.adsabs.harvard.edu/abs/2024MNRAS.528.6680Y}{Simulating image coaddition with the Nancy Grace Roman Space Telescope - II. Analysis of the simulated images and implications for weak lensing},
\textit{Monthly Notices of the Royal Astronomical Society}, Volume 528, Issue 4, March 2024, Pages 6680--6705, \Colorhref{https://doi.org/10.1093/mnras/stae177}{doi:10.1093/mnras/stae177}}
% \newline \textit{(``We have added \textbf{Kaili Cao} to the author list because of his contributions to finding some of the \textquotesingle{}known issues\textquotesingle{} and addressing them in the current working version of the code.'')}}

\cvitem{Feb 2024}{Christopher M. Hirata, Masaya Yamamoto, Katherine Laliotis, et al. (with \textbf{Kaili Cao} as 7th of 19 authors), \Colorhref{https://ui.adsabs.harvard.edu/abs/2024MNRAS.528.2533H}{Simulating image coaddition with the Nancy Grace Roman Space Telescope - I. Simulation methodology and general results},
\textit{Monthly Notices of the Royal Astronomical Society}, Volume 528, Issue 2, February 2024, Pages 2533--2561, \Colorhref{https://doi.org/10.1093/mnras/stae182}{doi:10.1093/mnras/stae182}}
% \newline \textit{(My contribution: Same as above.)}}

\cvitem{Sep 2021}{David J. Barnes, Mark Vogelsberger, Francesca A. Pearce, et al. (with \textbf{Kaili Cao} as 6th of 8 authors), \Colorhref{https://ui.adsabs.harvard.edu/abs/2021MNRAS.506.2533B}{Characterizing hydrostatic mass bias with MOCK-X},
\textit{Monthly Notices of the Royal Astronomical Society}, Volume 506, Issue 2, September 2021, Pages 2533--2550, \Colorhref{https://doi.org/10.1093/mnras/stab1276}{doi:10.1093/mnras/stab1276}}
% \newline \textit{(My contribution: I provided feedback on early data products of the \textsc{mock-X} framework, participated in discussions about relaxation criteria, and commented on the manuscript.)}}

\subsection{Co-authored Publications in Stellar Astrophysics}

\cvitem{Jul 2026}{Marc H. Pinsonneault, Jennifer L. van Saders, Lyra Cao, et al. (with \textbf{Kaili Cao} as 13th of 14 authors), \Colorhref{https://ui.adsabs.harvard.edu/abs/2026ApJ..1006..220P}{The YREC Stellar Evolution Code: Public Data Release},
\textit{The Astrophysical Journal}, Volume 1006, Number 2, July 2026, Page 220, \Colorhref{https://doi.org/10.3847/1538-4357/ae71cf}{doi:10.3847/1538-4357/ae71cf} [\Colorhref{https://github.com/yreclab/yrec.git}{GitHub}]}
% \newline \textit{(My contribution: I participated in the maintenance and modernization of the YREC code and related discussions and commented on the manuscript.)}}

\cvitem{Oct 2025}{B. Liagre, R.~A. Garc{\'\i}a, S. Mathur, et al. (with \textbf{K. Cao} as 7th of 12 authors), \Colorhref{https://ui.adsabs.harvard.edu/abs/2025A&A...702A.144L}{Beyond the Nyquist frequency: Asteroseismic catalog of undersampled Kepler late subgiants and early red giants},
\textit{Astronomy \& Astrophysics}, Volume 702, October 2025, Article Number A144, \Colorhref{https://doi.org/10.1051/0004-6361/202555167}{doi:10.1051/0004-6361/202555167}}
% \newline \textit{(My contribution: I provided theoretical predictions of first dredge-up signals and participated in related discussions.)}}

\cvitem{Jul 2025}{Trevor J. Weiss, Noah J. Downing, Marc H. Pinsonneault, et al. (with \textbf{Kaili Cao} as 7th of 13 authors), \Colorhref{https://ui.adsabs.harvard.edu/abs/2025ApJ...987..181W}{Modeling Asteroseismic Yields for the Roman Galactic Bulge Time-domain Survey},
\textit{The Astrophysical Journal}, Volume 987, Number 2, July 2025, Page 181, \Colorhref{https://doi.org/10.3847/1538-4357/adde5b}{doi:10.3847/1538-4357/adde5b}}
% \newline \textit{(My contribution: I participated in discussions about Roman asteroseismology and commented on the manuscript.)}}

\cvitem{May 2025}{M. Vrard, M.~H. Pinsonneault, Y. Elsworth, et al. (with \textbf{K. Cao} as 11th of 15 authors), \Colorhref{https://ui.adsabs.harvard.edu/abs/2025A&A...697A.165V}{Red giant evolutionary status determination: The complete Kepler catalog},
\textit{Astronomy \& Astrophysics}, Volume 697, May 2025, Article Number A165, \Colorhref{https://doi.org/10.1051/0004-6361/202452635}{doi:10.1051/0004-6361/202452635}}
% \newline \textit{(My contribution: I participated in discussions about the distinction between first-ascent red giant branch and asymptotic giant branch stars.)}}

\cvitem{Feb 2025}{Marc H. Pinsonneault, Joel C. Zinn, Jamie Tayar, et al. (with \textbf{Kaili Cao} as 19th of 41 authors), \Colorhref{https://ui.adsabs.harvard.edu/abs/2025ApJS..276...69P}{APOKASC-3: The Third Joint Spectroscopic and Asteroseismic Catalog for Evolved Stars in the Kepler Fields},
\textit{The Astrophysical Journal Supplement Series}, Volume 276, Number 2, February 2025, Page 69, \Colorhref{https://doi.org/10.3847/1538-4365/ad9fef}{doi:10.3847/1538-4365/ad9fef}}
% \newline \textit{(My contribution: I participated in discussions about the APOKASC-3 catalog, with significant contributions to uncertainties for spectroscopic observables, and contributed to paper writing.)}}

\cvitem{May 2024}{John D. Roberts, Marc H. Pinsonneault, Jennifer A. Johnson, et al. (with \textbf{Kaili Cao} as 11th of 19 authors), \Colorhref{https://ui.adsabs.harvard.edu/abs/2024MNRAS.530..149R}{Nature versus nurture: distinguishing effects from stellar processing and chemical evolution on carbon and nitrogen in red giant stars},
\textit{Monthly Notices of the Royal Astronomical Society}, Volume 530, Issue 1, May 2024, Pages 149--166, \Colorhref{https://doi.org/10.1093/mnras/stae820}{doi:10.1093/mnras/stae820}}
% \newline \textit{(My contribution: I contributed to uncertainties for spectroscopic observables in the APOKASC-3 catalog, provided theoretical insights into the first dredge-up, and commented on the manuscript.)}}

\fi

\section{Selected Presentations}

\subsection{2026}

\full{\cvitem{Jul 2026}{\textbf{Fisher Forecasts for Cosmological Yields from 3$\times$2pt Analysis of the Roman Space Telescope High Latitude Imaging Survey (poster)}, \textsl{Shifting Landscapes in Astrophysics: New Frontiers to Explore with Roman}, Pasadena, United States. [\Colorhref{https://workshop.ipac.caltech.edu/roman2026/files/KAILI_CAO_roman2026a.pdf}{Poster}, \Colorhref{https://www.youtube.com/watch?v=rXuoYL4ZUhQ\&t=115s}{Lightning talk video}, \Colorhref{https://kailicao.github.io/about/presentations/2026-07-13_roman-2026-lightning-talk.pdf}{Lightning talk slides}]}}

\full{\cvitem{Jul 2026}{\textbf{Effortless: Efficient Optimal Image Reconstruction for Roman and Beyond (poster)}, \textsl{Shifting Landscapes in Astrophysics: New Frontiers to Explore with Roman}, Pasadena, United States. [\Colorhref{https://workshop.ipac.caltech.edu/roman2026/files/KAILI_CAO_roman2026b.pdf}{Poster}, \Colorhref{https://www.youtube.com/watch?v=rXuoYL4ZUhQ\&t=115s}{Lightning talk video}, \Colorhref{https://kailicao.github.io/about/presentations/2026-07-13_roman-2026-lightning-talk.pdf}{Lightning talk slides}]}}

\cvitem{Jun 2026}{\textbf{Renovating Analysis Infrastructure for the Next Generation of Astronomical Surveys (dissertation defense)}, \textsl{The~Ohio~State~University}, Columbus, United States. [\Colorhref{https://kailicao.github.io/about/presentations/2026-06-25_dissertation-defense.pdf}{Slides}]}

\cvitem{Apr 2026}{\textbf{Weak Gravitational Lensing Cosmology with the Roman Space Telescope High Latitude Imaging Survey (seminar)}, \textsl{Cosmology Group Seminar (University~of~Michigan)}, Ann Arbor, United States. [\Colorhref{https://kailicao.github.io/about/presentations/2026-04-17_um-cosmology-seminar.pdf}{Slides}]}

\cvitem{Mar 2026}{\textbf{Hot-wiring the Roman Sky: Resolving the Undersampling Crisis with Fast IMCOM for Real-Time Discovery (contributed talk)}, \textsl{RAPID Response: Hot-wiring the Next Generation of Time-Domain Science (Caltech/IPAC)}, Pasadena, United States. (Fast IMCOM is the earlier name of Effortless.) [\Colorhref{https://kailicao.github.io/about/presentations/2026-03-25_rapid-response.pdf}{Slides}]}

\full{\cvitem{Jan 2026}{\textbf{Fisher Forecasts for Cosmological Yields from 3$\times$2pt Analysis of the Roman Space Telescope High Latitude Imaging Survey (poster)}, \textsl{Arizona Winter School on Multi-Probe Cosmology in the Roman ST Era (University of Arizona)}, Tucson, United States. [\Colorhref{https://kailicao.github.io/about/presentations/2026-01_multi-probe-winter-school-poster.pdf}{Poster}, an earlier version of the July 2026 poster]}}

\full{\cvitem{Jan 2026}{\textbf{Fast IMCOM: Unlocking New Scientific Frontiers with High-Fidelity Image Coaddition (contributed oral presentation)}, \textsl{247th Meeting of the American Astronomical Society}, Phoenix, United States. (Fast IMCOM is the earlier name of Effortless.) [\Colorhref{https://kailicao.github.io/about/presentations/2026-01-06_aas-247-fast-imcom.pdf}{Slides}, \Colorhref{https://ui.adsabs.harvard.edu/abs/2026AAS...24723508C/abstract}{Abstract}]}}

\subsection{2025}

\full{\cvitem{Nov 2025}{\textbf{Early Career Scientist Talk: Stage IV Observation, Stage IV Computation}, \textsl{Roman HLIS Cosmology PIT Face-to-Face Meeting 2025 (University of Pennsylvania)}, Philadelphia, United States. [\Colorhref{https://kailicao.github.io/about/presentations/2025-11-19_hlis-f2f-ecs-talk.pdf}{Slides}]}}

\cvitem{Oct 2025}{\textbf{Testing Stellar Theory with Joint Spectroscopic and Asteroseismic Catalogs (invited seminar talk)}, \textsl{AstroCoffee (University of Pittsburgh)}, Pittsburgh, United States. [\Colorhref{https://kailicao.github.io/about/presentations/2025-10-20_astrocoffee-pitt.pdf}{Slides}]}

\full{\cvitem{Oct 2025}{\textbf{Forecasts for weak lensing cosmology with Roman Space Telescope (contributed parallel talk)}, \textsl{28th International Conference on Particle Physics and Cosmology (COSMO-25)}, Pittsburgh, United States. [\Colorhref{https://kailicao.github.io/about/presentations/2025-10-13_cosmo-25.pdf}{Slides}]}}

\cvitem{Oct 2025}{\textbf{Toward Accurate Asteroseismic Masses for Luminous Giants (invited talk, online)}, \textsl{Galactic Genesis Working Group (SDSS-V Milky Way Mapper)}. [\Colorhref{https://kailicao.github.io/about/presentations/2025-10-07_galactic-genesis.pdf}{Slides}]}

\full{\cvitem{Aug 2025}{\textbf{PyIMCOM Hack Session 2025}, \textsl{Shear and Clustering Measurement Working Group Hack Week (Roman HLIS Cosmology PIT)}, Columbus, United States. [\Colorhref{https://kailicao.github.io/about/presentations/2025-08-04_pyimcom-hack-session.pdf}{Slides}]}}

\full{\subsection{2024 and Before}}

\full{\cvitem{Jul 2024}{\textbf{IMCOM: optimized image coaddition pipeline for Roman weak lensing cosmology (contributed talk)}, \textsl{Challenging Theory with Roman: From Planet Formation to Cosmology}, Pasadena, United States.
[\Colorhref{https://conference.ipac.caltech.edu/roman2024/abstracts}{Abstract} (please search for ``IMCOM''), \Colorhref{https://conference.ipac.caltech.edu/roman2024/talks/CAO_roman2024.pdf}{Slides}, \Colorhref{https://youtu.be/da4yKh0owPo}{Video}]}}

\full{\cvitem{Sep 2022}{\textbf{Theoretical predictions concerning the first dredge-up (presentation)}, \textsl{Summer Research Poster Session (The~Ohio~State~University)}, Columbus, United States. [\Colorhref{https://kailicao.github.io/about/presentations/2022-09_first-dredge-up.pdf}{Slides}]}}

\full{\cvitem{Aug 2019}{\textbf{How to define a relaxed cluster?}, \textsl{Summer MKI Undergraduate Research Forum (SMURF; Massachusetts Institute of Technology)}, Cambridge, United States. [\Colorhref{https://kailicao.github.io/about/presentations/2019-08-23_smurf-relaxed-clusters.pdf}{Slides}]}}

\section{Teaching Experience}

\full{\cventry{2025}{Microteaching Lesson, College Teaching (ESEPSY 7404)}{The~Ohio~State~University}{Columbus}{United States}{Planned a lesson by backward design, ``Basic Counting: Permutations and Combinations,'' and taught it to my peers, mostly from outside mathematics. Scored 33/35 by the instructor; median 48/50 from 14 classmates. [\Colorhref{https://kailicao.github.io/about/presentations/2025-04-01_microteaching-basic-counting.pdf}{Slides}]}}
\cventry{2022 -- 2023}{Instructor, Online Python Tutorial}{}{}{}{46 students in 4 cohorts\full{, mostly first-time programmers from the humanities and social sciences}; 28 took an active part. Gave lectures, designed exercises,\full{ revised the study guide from students' feedback,} and mentored personal projects.}
\cventry{2021 -- 2023}{Teaching Associate (Grader), Physics 7501 \& 7502: (Graduate) Quantum Mechanics}{The~Ohio~State~University}{Columbus}{United States}{}
\full{\cvitem{Miscellany}{Taught or tutored college mathematics (2017, 2018, 2020, 2023), college physics (2023), graduate school application (2022, 2023, 2024), French phonology (2017, 2018), etc.}}

\section{Service}

\cvitem{2026 -- present}{\textbf{Lead}, MapHub (the MIDAS arXiv Preprint Hub), a carpentry of the Schmidt AI in Science Fellowship, with Kyle Heyblom, Eva Albalghiti and Dimitrios Simatos.}
\cvitem{2025}{\textbf{Tutorial}, PyIMCOM, at an internal hack week of the Roman HLIS Cosmology Project Infrastructure Team.}

\section{Skills}

\cvitem{Programming}{Python (main), C/C++, Fortran, Shell/Perl, Wolfram Mathematica}
\cvitem{Languages}{Chinese (native), English (working proficiency), French (undergraduate major)}
\full{\cvitem{Certificates}{Productivity Program (The~Ohio~State~University, 2021), Responsible Conduct of Research (The~Ohio~State~University, 2024), Data Science Boot Camp (The Erd\H{o}s Institute, 2025)}}

\full{\section{References}

\cvitem{}{Available on request.}}

\end{document}
